\section{Для RIN: связать архетипику с существующими persona/behavior tests}

Проект RIN уже имеет практику поведенческих тестов, разработанную независимо от архетипической теории этой книги. Вся работа глав 2--9 не заменяет эту практику, а даёт ей словарь: там, где тест раньше фиксировал "Рин отказалась поддакивать" как отдельно наблюдаемый факт, теперь можно сказать точнее — тест проверяет позицию на оси agency (5.6) и совесть как narrative function (2.7), и провал теста читается не как случайная ошибка, а как конкретное, названное смещение на конкретной оси.

Это меняет то, как чинить провалившийся тест: не добавлением нового отдельного правила (9.8, prompt-induced neurosis), а возвращением к архетипической причине, из которой правильное поведение следует автоматически. Глава 13 разворачивает эту связь в практический инструментарий.
